% ==============================
% 4. Обзор существующих алгоритмов
% ==============================
\section{Обзор существующих алгоритмов}

\subsection{NP-трудность задачи размещения таблиц}
Задача размещения логических таблиц классификации по стадиям конвейера RMT является вычислительно сложной. В работе \cite{vass2020compiling} доказано, что эта задача является NP-трудной. Более того, авторы показывают, что для ряда моделей не существует алгоритма, который решал бы её точно за полиномиальное время. 

Следовательно, для практического применения при значительных объёмах входных данных (десятки и сотни таблиц) необходимо использовать приближённые методы, в частности, жадные эвристики.

\subsection{Точные методы: целочисленное линейное программирование (ILP)}
Точное решение задачи размещения может быть получено с помощью целочисленного линейного программирования (Integer Linear Programming, ILP). В работе \cite{jose2015compiling} предложена ILP-модель, включающая:
\begin{itemize}
	\item переменные, указывающие, в каких стадиях размещены таблицы;
	\item ограничения, гарантирующие соблюдение ёмкости памяти и зависимостей;
	\item целевую функцию, минимизирующую количество использованных стадий.
\end{itemize}
ILP позволяет найти оптимальное размещение, однако время его работы экспоненциально зависит от числа таблиц. Эксперименты из \cite{jose2015compiling} показывают, что для программы с 24 логическими таблицами  ILP-решатель затрачивает около 6300 секунд на поиск оптимума. Для более крупных программ время становится неприемлемым для практического использования, особенно при необходимости перекомпиляции в реальном времени.

\subsection{Жадные эвристики}
В связи с NP-трудностью задачи для практических применений разработаны жадные эвристики, которые работают за полиномиальное время и обеспечивают приемлемое качество размещения. В литературе \cite{jose2015compiling} выделяют четыре основных типа эвристик:

\paragraph{FFD (First Fit Decreasing)} – классический алгоритм упаковки. Таблицы сортируются по убыванию количества требуемых блоков памяти. Затем каждая таблица размещается в первую стадию, в которой для неё достаточно свободной памяти. FFD не учитывает топологический порядок зависимостей, что может приводить к невозможности разместить корректно все таблицы при наличии match-зависимостей.

\paragraph{FFL (First Fit by Level)} – перед сортировкой вычисляются уровни таблиц: уровень таблицы равен длине самого длинного пути в графе зависимостей от любого источника. Таблицы сортируются по возрастанию уровня, а внутри уровня – в произвольном порядке. Это гарантирует, что предки обрабатываются раньше потомков, однако из-за произвольного порядка внутри уровня это может приводить к неэффективному использованию памяти на одной стадии.

\paragraph{FFLS (First Fit by Level and Size)} – модификация FFL, в которой внутри одного уровня таблицы дополнительно сортируются по убыванию размера. Данный подход сочетает преимущества топологического порядка и плотной упаковки. 

\paragraph{MCF (Most Constrained First)} – таблицы сортируются по степени «ограниченности»: в первую очередь обрабатываются таблицы, которые могут быть размещены только в определённом типе памяти (например, тернарные – только в TCAM) или имеют наибольшую ширину ключа.


\subsection{Выводы по обзору}
Таким образом, задача размещения таблиц является NP-трудной. Точные методы (ILP) обеспечивают оптимальность, но неприменимы для больших программ из-за экспоненциального времени работы. Жадные эвристики работают быстрее. Однако классические жадные эвристики имеют недостаток – они могут не найти размещение для всех таблиц из-за однопроходности. Поэтому актуальной является разработка многопроходных и гибридных подходов. В последующих разделах будет представлена реализация многопроходного FFD и FFLS, а также результаты их экспериментального сравнения по времени работы и колличеству использованных стадий.